\documentclass[12pt, a4paper]{article}
\usepackage[utf8]{inputenc}
\usepackage[spanish]{babel}
\usepackage{amsmath}
\usepackage{graphicx}
\usepackage{hyperref}
\usepackage{listings}
\usepackage{color}
\usepackage{geometry}

\geometry{top=2.5cm, bottom=2.5cm, left=2.5cm, right=2.5cm}

\definecolor{codegreen}{rgb}{0,0.6,0}
\definecolor{codegray}{rgb}{0.5,0.5,0.5}
\definecolor{codepurple}{rgb}{0.58,0,0.82}
\definecolor{backcolour}{rgb}{0.95,0.95,0.92}

\lstdefinestyle{mystyle}{
    backgroundcolor=\color{backcolour},   
    commentstyle=\color{codegreen},
    keywordstyle=\color{magenta},
    numberstyle=\tiny\color{codegray},
    stringstyle=\color{codepurple},
    basicstyle=\ttfamily\footnotesize,
    breakatwhitespace=false,         
    breaklines=true,                 
    captionpos=b,                    
    keepspaces=true,                 
    numbers=left,                    
    numbersep=5pt,                  
    showspaces=false,                
    showstringspaces=false,
    showtabs=false,                  
    tabsize=2
}

\lstset{style=mystyle}

\title{Guía de Estudio Detallada: Proyecto de Sistemas Inteligentes\\
\large Análisis Teórico y Desglose de Código}
\author{Laboratorio \#3 - INFO1157}
\date{\today}

\begin{document}

\maketitle
\tableofcontents
\newpage

\section{Introducción}
Este documento proporciona una guía exhaustiva del proyecto, combinando conceptos teóricos con un análisis detallado del código fuente. Se explica cada archivo, clase y método implementado.

\section{Análisis de Código Fuente}

\subsection{1. Configuración Global (\texttt{config.py})}
Este archivo centraliza todas las constantes y rutas del proyecto para mantener el código ordenado y fácil de modificar.

\subsubsection{Rutas y Directorios}
Se definen las rutas base utilizando \texttt{pathlib} para asegurar compatibilidad entre sistemas operativos (Windows/Linux).
\begin{lstlisting}[language=Python]
import os
from pathlib import Path

# Rutas principales
PROJECT_ROOT = Path(__file__).parent.absolute()
DATA_DIR = PROJECT_ROOT / "data"
RESULTS_DIR = PROJECT_ROOT / "results"
NOTEBOOKS_DIR = PROJECT_ROOT / "notebooks"
\end{lstlisting}

\subsubsection{Configuración de Algoritmos}
Se definen diccionarios con los hiperparámetros por defecto para los algoritmos. Esto permite cambiar la configuración en un solo lugar sin tocar la lógica de los modelos.
\begin{lstlisting}[language=Python]
# Configuración de KMeans
KMEANS_CONFIG = {
    "n_clusters": 3,
    "random_state": 42,
    "n_init": 10,
    "max_iter": 300,
}

# Configuración de MeanShift
MEANSHIFT_CONFIG = {
    "bandwidth": "scott",
    "cluster_all": True,
    "quantile": 0.2,
}
\end{lstlisting}

\newpage
\subsection{2. Carga de Datos (\texttt{src/data\_loader.py})}
Este módulo se encarga de leer los archivos crudos (formato \texttt{.npy}) y entregarlos en formatos útiles (DataFrames o Arrays).

\subsubsection{Clase DataLoader}
Utiliza métodos estáticos (\texttt{@staticmethod}) para no necesitar instanciar la clase.

\textbf{Carga de Puntos 3D:}
Lee el archivo \texttt{puntos\_3d.npy} y lo convierte inmediatamente a un DataFrame de Pandas con columnas x, y, z para facilitar su visualización.
\begin{lstlisting}[language=Python]
    @staticmethod
    def load_3d_points():
        path = config.DATA_FILES["puntos_3d"]
        data = np.load(path)
        df = pd.DataFrame(data, columns=['x', 'y', 'z'])
        return df
\end{lstlisting}

\textbf{Información de Datos:}
Calcula estadísticas básicas de los datos cargados, como el número total de puntos y la distribución por clases.
\begin{lstlisting}[language=Python]
    @staticmethod
    def get_data_info(data, labels):
        unique_labels = np.unique(labels)
        info = {
            "total_points": len(data),
            "total_classes": len(unique_labels),
            "samples_per_class": {
                cls: np.sum(labels == cls) for cls in unique_labels
            }
        }
        return info
\end{lstlisting}

\newpage
\subsection{3. Modelo K-Means (\texttt{src/clustering/kmeans\_model.py})}
Implementa el algoritmo de agrupamiento K-Means.

\subsubsection{Inicialización (\texttt{\_\_init\_\_})}
El constructor acepta parámetros opcionales. Si no se entregan, toma los valores por defecto de \texttt{config.py}.
\begin{lstlisting}[language=Python]
class KMeansModel:
    def __init__(self, n_clusters=None, random_state=None, ...):
        # Usar valores de config si no se proporcionan
        n_clusters = n_clusters or config.KMEANS_CONFIG["n_clusters"]
        
        self.model = SklearnKMeans(
            n_clusters=n_clusters,
            ...
        )
\end{lstlisting}

\subsubsection{Entrenamiento (\texttt{fit})}
Ajusta el modelo a los datos y guarda los resultados importantes (etiquetas, centroides, inercia) como atributos de la clase para fácil acceso.
\begin{lstlisting}[language=Python]
    def fit(self, X):
        self.model.fit(X)
        self.labels = self.model.labels_
        self.centroids = self.model.cluster_centers_
        self.inertia = self.model.inertia_
        return self
\end{lstlisting}

\newpage
\subsection{4. Modelo MeanShift (\texttt{src/clustering/meanshift\_model.py})}
Implementa MeanShift, un algoritmo que no requiere saber el número de clusters de antemano.

\subsubsection{Estimación de Bandwidth}
Una parte crítica de MeanShift es el "ancho de banda". Si se configura como "scott", el método \texttt{fit} calcula automáticamente este valor basándose en la densidad de los datos.
\begin{lstlisting}[language=Python]
    def fit(self, X):
        # Estimar ancho de banda si es necesario
        if self.bandwidth == "scott":
            bandwidth = estimate_bandwidth(X, quantile=self.quantile)
        else:
            bandwidth = self.bandwidth
        
        self.model = SklearnMeanShift(bandwidth=bandwidth, ...)
        self.model.fit(X)
\end{lstlisting}

\newpage
\subsection{5. Regresión Lineal (\texttt{src/regression/linear\_regression.py})}
Modelo para predecir una variable continua (peso) a partir de otra (altura).

\subsubsection{Predicción y Validación}
El método \texttt{predict} verifica primero si el modelo ha sido entrenado (\texttt{\_is\_fitted}) para evitar errores.
\begin{lstlisting}[language=Python]
    def predict(self, X):
        if not self._is_fitted:
            raise ValueError("El modelo debe ser entrenado primero")
        
        self.predictions = self.model.predict(X)
        return self.predictions
\end{lstlisting}

\subsubsection{Obtención de la Ecuación}
Este método auxiliar devuelve la fórmula de la recta $y = mx + b$ en formato legible, útil para visualización.
\begin{lstlisting}[language=Python]
    def get_equation(self):
        m, b = self.get_coefficients()
        return f"Peso = {m:.4f} x Altura + {b:.4f}"
\end{lstlisting}

\newpage
\subsection{6. Similitud Coseno (\texttt{src/similarity/cosine\_distance.py})}
Calcula qué tan similares son dos vectores basándose en su ángulo.

\subsubsection{Cálculo Matemático}
Implementa la fórmula del coseno manualmente usando NumPy. Se añaden protecciones contra división por cero.
\begin{lstlisting}[language=Python]
    @staticmethod
    def cosine_distance(vector_a, vector_b):
        # Producto punto
        dot_product = np.dot(a, b)
        # Normas (magnitudes)
        norm_a = np.linalg.norm(a)
        norm_b = np.linalg.norm(b)
        
        if norm_a == 0 or norm_b == 0:
            return 0.0
            
        return dot_product / (norm_a * norm_b)
\end{lstlisting}

\subsubsection{Conversión de Ángulos}
Métodos auxiliares para interpretar el resultado en grados, que es más intuitivo para humanos que radianes o el valor del coseno puro.
\begin{lstlisting}[language=Python]
    @staticmethod
    def angle_from_cosine_degrees(cos_theta):
        # Asegura que el valor esté en [-1, 1]
        cos_theta = np.clip(cos_theta, -1.0, 1.0)
        radians = np.arccos(cos_theta)
        return np.degrees(radians)
\end{lstlisting}

\newpage
\subsection{7. Filtro de Kalman (\texttt{src/filters/kalman\_filter.py})}
Estima el estado de un sistema dinámico lineal.

\subsubsection{Inicialización de Matrices}
Define las matrices clave del filtro: Transición ($F$), Medición ($H$), Covarianza ($P$) y Ruido del Proceso ($Q$).
\begin{lstlisting}[language=Python]
    def __init__(self, dt=0.1, ...):
        # Matriz de transición de estado (F)
        # x = x + v*dt
        # v = v
        self.F = np.array([[1, dt], [0, 1]])
        
        # Matriz de medición (H)
        # Solo medimos posición (1, 0)
        self.H = np.array([[1, 0]])
\end{lstlisting}

\subsubsection{Ciclo Predecir-Actualizar}
El método \texttt{filter} itera sobre una lista de mediciones, ejecutando el ciclo de Kalman en cada paso.
\begin{lstlisting}[language=Python]
    def filter(self, measurements):
        estimates = []
        for z in measurements:
            self.predict()  # Proyectar estado al futuro
            self.update(z)  # Corregir con nueva medición
            estimates.append(self.x[0, 0])
        return np.array(estimates)
\end{lstlisting}

\newpage
\subsection{8. Filtro Complementario (\texttt{src/filters/complementary\_filter.py})}
Fusiona señales de alta y baja frecuencia.

\subsubsection{Fusión de Señales}
Aplica la fórmula clásica del filtro complementario. $\alpha$ (alpha) controla cuánto confiamos en cada sensor.
\begin{lstlisting}[language=Python]
    def filter(self, measurement_noisy, measurement_integrated):
        # alpha * (dato ruidoso) + (1-alpha) * (dato integrado)
        estimate = (self.alpha * measurement_noisy + 
                   (1 - self.alpha) * measurement_integrated)
        
        self.position = estimate
        return estimate
\end{lstlisting}

\subsubsection{Búsqueda de Alpha Óptimo}
Un método estático que prueba múltiples valores de $\alpha$ para encontrar cuál minimiza el error respecto a una referencia conocida.
\begin{lstlisting}[language=Python]
    @staticmethod
    def find_optimal_alpha(noisy, integrated, reference):
        alphas = np.linspace(0, 1, 21)
        errors = []
        
        for alpha in alphas:
            cf = ComplementaryFilter(alpha=alpha)
            # ... calcular error ...
            errors.append(error)
        
        return alphas[np.argmin(errors)]
\end{lstlisting}

\newpage
\section{Notebooks de Jupyter}
Esta sección detalla el contenido y propósito de los notebooks interactivos utilizados para ejecutar y visualizar los algoritmos.

\subsection{Notebook A: Clustering y Regresión (\texttt{notebooks/parte\_a.ipynb})}
Este notebook cubre la primera parte del laboratorio, enfocándose en algoritmos de agrupamiento y regresión lineal.

\subsubsection{Estructura del Notebook}
\begin{itemize}
    \item \textbf{Punto 1: KMeans Clustering}
    \begin{itemize}
        \item Carga los datos de entrenamiento (\texttt{A.npy}).
        \item Aplica el modelo \texttt{KMeansModel} para encontrar 3 clusters.
        \item Visualiza los datos originales y los clusters resultantes con sus centroides.
    \end{itemize}
    \item \textbf{Punto 2: MeanShift Clustering}
    \begin{itemize}
        \item Carga los datos (\texttt{X.npy}).
        \item Utiliza \texttt{MeanShiftModel} para detectar automáticamente el número de clusters.
        \item Compara los resultados con KMeans.
    \end{itemize}
    \item \textbf{Punto 3: Regresión Lineal}
    \begin{itemize}
        \item Entrena un modelo para predecir pesos basados en alturas.
        \item Calcula el error cuadrático medio (MSE) y visualiza la recta de regresión.
    \end{itemize}
    \item \textbf{Punto 4: Distancia del Coseno}
    \begin{itemize}
        \item Calcula la similitud entre varios vectores de prueba.
        \item Muestra los ángulos resultantes para interpretar la similitud.
    \end{itemize}
\end{itemize}

\subsubsection{Código Clave (KMeans)}
\begin{lstlisting}[language=Python]
# Cargar datos
A, labels = DataLoader.load_kmeans_data()

# Inicializar y entrenar modelo
kmeans = KMeansModel(n_clusters=3)
kmeans.fit(A)

# Obtener centroides
centroids = kmeans.get_centroids()
\end{lstlisting}

\subsection{Notebook B: Visualización 3D y Filtros (\texttt{notebooks/parte\_b.ipynb})}
Este notebook aborda la visualización avanzada y los algoritmos de filtrado.

\subsubsection{Estructura del Notebook}
\begin{itemize}
    \item \textbf{Punto 5: Visualización 3D}
    \begin{itemize}
        \item Carga una nube de puntos 3D (\texttt{puntos\_3d.npy}).
        \item Calcula el centroide geométrico.
        \item Genera un gráfico interactivo 3D mostrando los puntos y el centroide.
    \end{itemize}
    \item \textbf{Punto 6: Dataset Iris}
    \begin{itemize}
        \item Carga el dataset Iris.
        \item Aplica clustering y evalúa el rendimiento con métricas ARI y Silhouette.
    \end{itemize}
    \item \textbf{Punto 7: Filtro de Kalman}
    \begin{itemize}
        \item Simula una trayectoria con ruido.
        \item Aplica el filtro de Kalman para suavizar la trayectoria y estimar la posición real.
    \end{itemize}
    \item \textbf{Punto 8: Filtro Complementario}
    \begin{itemize}
        \item Simula datos de sensores (acelerómetro y giroscopio).
        \item Fusiona las señales para obtener una estimación de ángulo más estable.
    \end{itemize}
\end{itemize}

\subsubsection{Código Clave (Visualización 3D)}
\begin{lstlisting}[language=Python]
# Cargar puntos
puntos_3d = np.load('../puntos_3d.npy')

# Calcular centroide
centroid = np.mean(puntos_3d, axis=0)

# Graficar
fig = plt.figure()
ax = fig.add_subplot(111, projection='3d')
ax.scatter(puntos_3d[:,0], puntos_3d[:,1], puntos_3d[:,2])
ax.scatter(centroid[0], centroid[1], centroid[2], c='red', s=100)
\end{lstlisting}

\newpage
\section{Apéndice: Código Fuente Completo}
A continuación se presenta el código fuente completo de todos los módulos implementados en el proyecto, con una breve descripción de su propósito.

\subsection{Configuración (config.py)}
Define las constantes globales, rutas de archivos y parámetros de configuración por defecto para todos los algoritmos.
\lstinputlisting[language=Python]{src/../config.py}

\subsection{Cargador de Datos (src/data\_loader.py)}
Clase utilitaria encargada de cargar archivos \texttt{.npy} y transformarlos en estructuras de datos manejables como DataFrames de Pandas.
\lstinputlisting[language=Python]{src/data_loader.py}

\subsection{K-Means Model (src/clustering/kmeans\_model.py)}
Implementación orientada a objetos del algoritmo K-Means, encapsulando la inicialización, entrenamiento y predicción de clusters.
\lstinputlisting[language=Python]{src/clustering/kmeans_model.py}

\subsection{MeanShift Model (src/clustering/meanshift\_model.py)}
Implementación del algoritmo MeanShift que permite la detección automática de clusters basada en densidad, incluyendo la estimación de ancho de banda.
\lstinputlisting[language=Python]{src/clustering/meanshift_model.py}

\subsection{Regresión Lineal (src/regression/linear\_regression.py)}
Modelo de regresión lineal simple para predecir variables continuas, incluyendo métodos para evaluación de métricas (MSE, R2).
\lstinputlisting[language=Python]{src/regression/linear_regression.py}

\subsection{Distancia Coseno (src/similarity/cosine\_distance.py)}
Módulo con métodos estáticos para calcular la similitud coseno entre vectores y convertir los resultados a ángulos en grados.
\lstinputlisting[language=Python]{src/similarity/cosine_distance.py}

\subsection{Filtro de Kalman (src/filters/kalman\_filter.py)}
Implementación de un filtro de Kalman unidimensional estándar, con métodos separados para las fases de predicción y actualización.
\lstinputlisting[language=Python]{src/filters/kalman_filter.py}

\subsection{Filtro Complementario (src/filters/complementary\_filter.py)}
Implementación de un filtro complementario para fusión de sensores, útil para combinar señales de alta y baja frecuencia (ej. acelerómetro y giroscopio).
\lstinputlisting[language=Python]{src/filters/complementary_filter.py}

\end{document}
